%% Response to reviewers, SEGAN-D-26-03850
%% Compile with pdfLaTeX.
\documentclass[11pt,a4paper]{article}
\usepackage[margin=2.4cm]{geometry}
\usepackage{amsmath,amssymb}
\usepackage[hidelinks]{hyperref}
\usepackage{booktabs}
\usepackage{xcolor}
\usepackage{mdframed}
\usepackage{parskip}
\usepackage{enumitem}
\usepackage{titlesec}

\definecolor{quotebg}{gray}{0.96}
\definecolor{quoterule}{gray}{0.75}
\newenvironment{revcomment}
  {\begin{mdframed}[backgroundcolor=quotebg,linecolor=quoterule,linewidth=0.4pt,
    innertopmargin=6pt,innerbottommargin=6pt,skipabove=8pt,skipbelow=8pt]\itshape}
  {\end{mdframed}}
\newcommand{\action}[1]{\textbf{Action.} #1}
\newcommand{\where}[1]{\textbf{Where.} #1}
\titleformat{\section}{\large\bfseries}{\thesection.}{0.6em}{}
\titleformat{\subsection}{\bfseries}{\thesubsection}{0.6em}{}

\title{\vspace{-1.4cm}\textbf{Response to the reviewers}\\[4pt]
\large Manuscript SEGAN-D-26-03850, \emph{Sustainable Energy, Grids and Networks}}
\author{Pablo Reyes Cerda \and Kerven Cea Morales}
\date{}

\begin{document}
\maketitle
\vspace{-1.0cm}

\section*{Letter to the Editor}

Dear Professor Chicco,

Thank you for the opportunity to revise this manuscript, and please pass our thanks to both reviewers. The reports were unusually specific, and one comment in particular changed the paper substantially.

Reviewer~1's first comment asked us to implement a competitive alternative base model, a multi-quantile gradient-boosting predictive, and to repeat every Transport+ACI experiment on it, in order to determine whether the advantage of our calibration scheme came from the calibration mechanism or from its compatibility with the specific base distribution. We did exactly that, and the answer is that it came from the base model. The gain we reported does not replicate. With the multi-quantile predictive, the static conformal split is already at nominal coverage in the window of maximum divergence and the adaptive scheme returns intervals that are marginally wider. Across three base models and five evaluation windows the width recovered is linearly related to how far the static split over-covers ($r=-0.755$), so what the adaptive layer was recovering is base-model over-dispersion, not distribution shift. The ablations Reviewer~1 requested in comment~3 localise the effect further: online adaptation is the component that works, score transport significantly worsens the interval score once adaptation is present, and the bootstrap tail shrinkage improves it by six tenths of one per cent while consuming roughly ninety per cent of the runtime.

We have not tried to save the original claim. The paper is reframed around the finding: the contribution is now the diagnostic itself, that the apparent benefit of an adaptive conformal layer measures the miscalibration of the base predictive it corrects, together with the experiment that establishes it and the practical ordering that follows, which is that fixing the base model dominates fixing the calibration layer. We also report, because it is the honest comparison, that a one-sided conformalized quantile regression on the better base model attains 89.8 per cent coverage at 276~MWh against 596~MWh for the pipeline we had proposed. This is a less flattering paper than the one we submitted. We think it is a more useful one, and we note that it directly addresses the concern behind Reviewer~1's answer to structured question~5, that the interpretation and conclusions were not supported by the data.

The dataset contribution is unchanged, and every number in it has been re-verified.

One matter concerns Reviewer~2's second comment, on the evidence for a storage-driven regime change. The wording the reviewer asked us to adopt is weaker than the position the manuscript already took: the submitted version made no causal attribution to storage, dated the change by a formal change-point analysis rather than by any technology's entry, and defined the evaluation window as the one of maximum divergence from the calibration distribution. We should nonetheless have been clearer, and the reviewer's comment let us find why the text read otherwise. A key in our code, \texttt{test\_ramp}, and its accompanying footnote mentioning storage deployment had survived from an earlier version of the analysis into the published table. That is a nomenclature residue, and it was reasonable to read a causal thesis into it. We have renamed the key to \texttt{test\_transition} throughout the code, the result files, the tables and the figures, removed the footnote, and audited every remaining mention of storage. We verified that the rename changed labels only: the regenerated result files are identical to the submitted ones in every numerical cell.

A third point of disclosure concerns how much of the advantage of the multi-quantile predictive belongs to its form. Its CRPS is eighteen to thirty per cent lower than the hurdle's across the evaluation windows, but the two arms share the training protocol and not the fitted budget: that advantage requires a budget of the order of five thousand trees, and at the hurdle's own budget of eight hundred trees it nearly disappears. Section~5.4 separates model form from fitted budget and reports both effects, including the part that reduces our claim. Section~5.4 also reports that the hurdle used throughout the paper is not the best available on this problem: refitted with early stopping, it selects an eighth of the trees and beats our base model by 5.7 per cent of CRPS. Section~7 lists every table and figure whose levels should be read with that discount. The discount runs in one direction and changes no ordering we report, including that of point forecasting, where the stopped hurdle is worse by mean absolute error, as Section~5.1 measures.

Two further points of disclosure. First, the hyperparameters $\gamma$ and the window length were not selected on real data in the submitted version; we have redone the selection by rolling origin inside the calibration window, as Reviewer~1 asked in comment~4, and it returns values an order of magnitude gentler than those we reported. Second, when we regenerated every descriptive claim of Section~3 from versioned scripts, six statements did not reproduce. None affects a result; all were descriptions of the dataset. They are corrected in the revised text and listed in Section~\ref{sec:correcciones} of this letter.

We respond to each of the fifteen comments below. Reviewer comments are quoted verbatim in grey, followed by the action taken and its exact location in the revised manuscript.

\bigskip
Yours sincerely,\\
Pablo Reyes Cerda and Kerven Cea Morales

\clearpage
\section*{Summary of changes}

\begin{tabular}{@{}p{0.13\textwidth}p{0.46\textwidth}p{0.35\textwidth}@{}}
\toprule
Comment & Change & Location \\
\midrule
R1.1 & Multi-quantile GBM base model implemented; all experiments repeated on eight base models spanning a capacity ladder; gain does not replicate; model form separated from fitted budget; paper reframed & \S4.1, \S5.3, \S5.4, Tables 4--7, Figs.~7--8, abstract, \S8 \\
R1.2 & Algorithm in full, flowchart, complete hyperparameter table with seeds, execution order stated & \S4.5, Alg.~1, Fig.~4, App.~C \\
R1.3 & Five ablation groups with marginal contribution, bootstrap intervals and runtime & \S5.5, Table~8, Fig.~9 \\
R1.4 & Rolling-origin selection inside the calibration window; test reserved for one evaluation; sensitivity reported as display & \S4.7, \S5.8 \\
R1.5 & One-sided interval score adopted; infinite limits handled explicitly; $\alpha$ bound evaluated; four benchmarks added; all differences with day-block bootstrap & \S4.8, \S5.6, Tables 9--10 \\
R2.1 & Bibliography expanded from 10 to 43 entries, every DOI verified, no grouped citations & \S2 \\
R2.2 & \texttt{test\_ramp} renamed, footnote removed, storage mentions audited; rename verified to change labels only & \S4.8, \S5.2, Table~3, Figs.~5--6, code \\
R2.3 & Single chronology; 180 detection configurations; six alternative window boundaries evaluated & \S3.4, \S5.8, Table~12 \\
R2.4 & Explicit statement of what is new and what is taken from the literature & \S4.6 \\
R2.5 & ICC and design effect quantified; block, group and plant-level conformal tested; coverage disaggregated on four axes & \S4.3, \S5.7, Table~11 \\
R2.6 & Empirical quantiles, CQR, rolling historical distributions and zero-inflated parametric benchmarks; six metrics reported & \S5.6, Table~9 \\
R2.7 & Point-forecasting claim scoped; three guarantee regimes separated in abstract, \S4.3 and conclusion & abstract, \S4.3, \S5.1, \S8 \\
Q4 & Four new figures & Figs.~4, 7, 8 and 9 \\
Q7 & Dedicated limitations section with ten items & \S7 \\
Q8, Q9 & Restructuring and language editing throughout & whole manuscript \\
Editorial & Title shortened from 26 to 14 words, country removed, system moved to abstract & title, abstract \\
\bottomrule
\end{tabular}

\clearpage
\section{Editorial note}

\begin{revcomment}
Editorial note: the title of the article is very long. Is it necessary to mention ``for Chile's National Electricity System'' in the title? Journal articles should have more general validity that the application to a single country.
\end{revcomment}

\action{Agreed on both counts. The title went from 26 words to 14 and no longer names the country. It keeps the two contributions in view, the dataset and the diagnostic, and it keeps the paper identifiable as a power-systems paper: an earlier draft of this revision named only the method, and on rereading we judged that a title for this journal should say what the work is about.}

Old: \emph{Probabilistic forecasting of renewable curtailment under regime change: an open dataset and a conformal method with distribution-shift-aware coverage control for Chile's National Electricity System.}

New: \emph{An open curtailment dataset and a diagnostic for adaptive conformal calibration under regime change.}

We considered two alternatives and record them here in case the Editor prefers one:

\begin{enumerate}[leftmargin=1.6em,itemsep=2pt]
\item \emph{Adaptive conformal calibration under regime change: the gain measures base-model miscalibration} --- states the finding directly, at the cost of naming neither the dataset nor the application domain.
\item \emph{When does adaptive conformal calibration pay? Evidence from renewable curtailment} --- foregrounds the question rather than the answer.
\end{enumerate}

\where{Title page. The description of the system, its size and the 19{,}000~GWh curtailed over 2022 to mid-2026, has moved into the abstract, second sentence, so that no information is lost. The abstract was also brought within the 250-word limit of the guide for authors, from 431 words to 248, without dropping the diagnostic result, the fitted relationship or the separation of the three guarantee regimes. The country still appears in the abstract, in Section~3 and in the keywords, where it is a statement of provenance rather than a limitation of scope.}

\clearpage
\section{Reviewer 1}

\subsection*{R1.1 --- Model-agnosticism of the calibration layer}

\begin{revcomment}
This work uses only the hurdle mixture model as the base probabilistic forecasting model, which is insufficient to validate the stability of the proposed calibration method across different base models. Although the hurdle model aims to fit a full predictive distribution instead of optimising point prediction errors, PIT scoring, transport mapping and ACI updates all rely on this base distribution. The quality of the base model directly determines the sharpness and calibration performance of final prediction intervals. Current results cannot tell whether the performance advantages of Transport+ACI originate from the calibration mechanism itself or merely its compatibility with this specific base distribution setup. It is recommended to adopt at least one competitive probabilistic forecasting model (e.g., multi quantile GBM), repeat all Transport+ACI experiments, and verify the model agnostic property and stability of the proposed method.
\end{revcomment}

This comment was correct, and answering it changed the paper.

\action{We implemented the multi-quantile gradient-boosting predictive the reviewer suggested and repeated every experiment on it. To separate two competing explanations we added a third arm, the hurdle with a conditional dispersion $\sigma(x)$, so that the comparison distinguishes the effect of constant dispersion from the effect of the lognormal form itself.}

The multi-quantile model estimates 54 quantiles from 0.005 to 0.999 with the pinball loss, under exactly the same twelve features, the same rows, the same 31 December 2023 training cut-off, the same a-priori hyperparameters and the same seed as the hurdle. Crossings are removed by rearrangement. Its zero inflation is not imposed: it is read off the grid as the highest level whose quantile is zero.

We refactored the calibration layer so that it touches the base model through exactly two operations, the predictive CDF and its generalised inverse. That makes the comparison exact rather than approximate: nothing else changes between arms, neither the score, nor the window boundaries, nor the seven-day embargo, nor the seeds, nor $\gamma$, nor the 60-day recent window, nor the shrinkage constants. As a control we assert inside the experiment that the hurdle arm reproduces Table~2 of the submitted version in all thirty cells and four metrics. It does.

\textbf{Result.} The advantage does not replicate. In the transition window:

\begin{center}
\begin{tabular}{@{}lccc@{}}
\toprule
Base model & Static split & Transport+ACI ($\gamma=0.05$) & Width change \\
\midrule
Hurdle, constant $\sigma$ & 95.3\% / 1203~MWh & 90.9\% / 806~MWh & $-33.1$\% \\
Hurdle, $\sigma(x)$ & 91.7\% / 770~MWh & 89.5\% / 733~MWh & $-4.8$\% \\
Multi-quantile GBM & 90.6\% / 410~MWh & 90.5\% / 425~MWh & $+3.6$\% \\
\bottomrule
\end{tabular}
\end{center}

With the multi-quantile predictive the width difference is $+15.1$~MWh with a day-block bootstrap interval of $[+0.0,+30.3]$, which only just excludes zero; the interval-score difference, $+13.9$ with an interval of $[-10.7,+37.8]$, contains it. On either reading the point estimate is on the opposite side from the one we reported. Over the fifteen cells given by three base models and five windows, the width recovered by the adaptive layer is linearly related to how far the static split over-covers: Pearson $r=-0.755$, slope $-4.36$ per cent of width per percentage point of over-coverage, intercept $+5.45$ per cent. When the static split is already at nominal, the adaptive layer costs width.

The reviewer's diagnosis was therefore right, and stronger than stated: not only can the reported results not distinguish the two explanations, the correct explanation is the second one. What the layer recovers is base-model over-dispersion.

We also record that the multi-quantile predictive dominates the hurdle on every proper metric and in every window, with a CRPS 18 to 30 per cent lower, and that in three of the five windows the best of the eighteen model-by-method combinations by interval score is the multi-quantile predictive with the plain static split and no adaptive layer at all.

That last point needed a control, and we have run it rather than flagging it. The arms share features, rows, training cut-off, per-model hyperparameters and seed, but they do not share capacity: the hurdle is two boosted models and the multi-quantile predictive is fifty-four. A reader could reasonably object that the comparison confounds the model form with the fitted budget. Section~5.4 separates the two by varying the tree budget at a fixed grid, and in the other direction by giving the hurdle five times its own.

We state the outcome plainly, including the part that goes against the previous wording. \textbf{At a strictly matched budget the advantage nearly disappears}: 810 trees against 800 improve the CRPS of the whole test by 2.4 per cent, not by the eighteen to thirty per cent we reported. That figure requires roughly five thousand trees, and the previous version of this response asserted it alongside a claim of identical discipline, which mixed the two effects. That was our error and it is corrected. What survives, and is a stronger result than the limitation it replaces, is that the hurdle cannot buy parity: given four thousand trees it gets \emph{worse}, its CRPS rising from 89.5 to 95.5~MWh, while at a comparable budget of 2{,}700 trees the multi-quantile predictive reaches 100.6 against 141.1. On this panel and between these two families, the multi-quantile form converts capacity into predictive accuracy and the two-part lognormal form does not. We claim nothing beyond that.

The mechanism corroborates a result already in the manuscript. When the hurdle is enlarged its estimated dispersion falls, from 1.7016 to 1.6641, and its intervals nonetheless widen by 292~MWh. The cause is the conformal quantile, which moves from 0.9242 to 0.9532; crossing the two fitted components shows that substituting only the enlarged $\mu(x)$ accounts for 93 per cent of that displacement and substituting only the smaller $\hat\sigma$ for almost none. The degradation is of location, not of dispersion, which is the experimental counterpart of what we already report for the conditional-dispersion variant.

\textbf{The comment also let us answer the stability question on more evidence than we first used.} Each arm of the ladder passes through the same calibration layer in the same five windows, so the diagnostic relationship can be estimated over forty cells from eight base models rather than fifteen from three. That matters because the fifteen cells span over-coverage only over $[-2.06, +5.34]$ percentage points, while the low-capacity arms reach $-8.42$ and populate the region below nominal where the original design had no point. The relationship strengthens: $r$ goes from $-0.755$ to $-0.873$ with a bootstrap interval of $[-0.912, -0.645]$, the slope from $-4.36$ to $-4.89$, and the intercept, which on fifteen cells was $[+0.07, +11.25]$ and barely excluded zero, becomes $[+1.60, +9.84]$. Leaving out any single cell moves $r$ within $[-0.889, -0.820]$ and leaving out any entire base model within $[-0.899, -0.792]$.

We considered stating the relationship symmetrically, since the low-capacity arms under-cover and the layer widens their intervals, which is the mirror of the over-covered case. We do not, because the data do not support it: fitting the two sides separately improves the fit by more than chance ($p=0.0032$), and the improvement is one of level rather than of rate ($p=0.0009$ for the level, $p=0.41$ for the slope). The claim therefore stays scoped to over-coverage, and the under-covered arms are reported as evidence that the mechanism does not stop at the nominal level rather than as a symmetric law.

\where{Section~4.1 describes the base models. Section~5.3 is a new subsection reporting the result, with Table~4 (the three arms in the transition window), Table~5 (the fifteen cells and the fitted relationship), Table~6 (the forty cells with their bootstrap intervals) and Figure~7. Section~5.4 is a further new subsection with the capacity ladder, Table~7 and Figure~8. The abstract, the introduction (contribution paragraphs), Section~6.1 and the conclusion are rewritten around it. Section~7 states the limits of the diagnostic. The commands that regenerate it are in Appendix~D.}

\subsection*{R1.2 --- Implementation detail and algorithm flowchart}

\begin{revcomment}
The description of the Transport+ACI method remains conceptual and lacks essential implementation details. Critical practical information is not adequately specified, including the length of the recent observation window, selection rules for empirical quantiles, interpolation schemes, bootstrap resampling repetitions, the exact tail shrinkage function, shrinkage strength, and the execution order between transport mapping computation and ACI updates. It is recommended that the authors add an algorithm flowchart to fully illustrate the complete workflow of Transport+ACI: from input processing, PIT score calculation, transport map estimation, tail robustification and ACI updates, to the generation of final prediction intervals.
\end{revcomment}

\action{All seven items are now specified in the body, and the flowchart has been added. We also added a complete hyperparameter table which is \emph{generated from the source code} rather than transcribed, by a script that imports the modules and reads the constants and default arguments, so that it cannot drift from the implementation.}

Taking the seven items in the reviewer's order:

\begin{itemize}[leftmargin=1.4em,itemsep=3pt]
\item \textbf{Length of the recent window.} 60 days in the reported configuration, ending at $t-7$ rather than at $t$ because of the horizon embargo. The window length is one of the two hyperparameters now selected by rolling origin (R1.4), which returns 120 days.
\item \textbf{Selection rule for empirical quantiles.} An equispaced grid of $g$ levels on $[0,1]$ with $g=\min(100,\max(20,m/2))$, where $m$ is the number of scores in the recent window; the second term prevents a grid finer than the data support.
\item \textbf{Interpolation scheme.} Two successive piecewise-linear interpolations, first from a score to its level on the old empirical quantile function, then from that level to the regularised new quantile. This is the increasing rearrangement evaluated pointwise.
\item \textbf{Bootstrap repetitions.} $B=150$ resamples of the recent window. The map is the mean of the replicate quantile functions, a bagged estimate; the replicate standard deviation at each level is the band.
\item \textbf{Exact tail shrinkage function.} Now given as equation~(5) of the manuscript: $\hat q_{\mathrm{reg}}(u)=(1-\lambda(u))\hat q_{\mathrm{old}}(u)+\lambda(u)\hat q^{\mathrm{bag}}_{\mathrm{new}}(u)$ with $\lambda(u)=[1+(\mathrm{sd}_{\mathrm{boot}}(u)/\tau)^2]^{-1}$, followed by a running maximum and a clip to $[0,1]$.
\item \textbf{Shrinkage strength.} $\tau=0.05$ in PIT units, fixed a priori and never tuned against the test set.
\item \textbf{Execution order.} Stated explicitly and marked in the figure. On date $t$: (i) if a refresh is due, recompute the map from the window ending at $t-7$ and transport the whole calibration pool; (ii) take the conformal quantile from the transported pool at the \emph{current} $\alpha_t$, which does not yet include the error of day $t$; (iii) emit the interval; (iv) push the error of day $t$ onto the embargo queue, which updates $\alpha$ seven steps later. The map never uses the adaptive level and the adaptive update never uses the map; the only coupling is the pool.
\end{itemize}

\where{Section~4.5 gives the three items above that belong in the body. Algorithm~1 gives the full procedure in pseudocode. Figure~4 is the flowchart the reviewer asked for, with the execution order marked along the top and the two branches drawn separately so the coupling point is visible. Appendix~C, Table~C.13, is the complete hyperparameter specification including all four seeds. Appendix~D gives the environment and the command sequence.}

\subsection*{R1.3 --- Ablations}

\begin{revcomment}
Ablation experiments for the Transport+ACI method are still inadequate, and the actual contribution of each core component to performance improvement cannot be clearly identified. Existing experiments lack targeted validation for the internal key mechanisms of Transport+ACI. In particular, the authors introduce a bootstrap variance based tail shrinkage strategy built upon score transport, yet experimental results with this strategy disabled are not reported. Therefore, it is difficult to judge whether performance gains are mainly attributed to score transport, ACI adaptive updates, or the tail robustification mechanism. The authors are suggested to add ablation groups including pure ACI, Transport+ACI (without tail shrinkage), and the full Transport+ACI model, to further verify whether newly introduced mechanisms and their computational complexity bring tangible performance improvements.
\end{revcomment}

\action{We ran the three groups requested and added two more to close the factorial, so that each component can be attributed unambiguously: static split, pure ACI, transport alone at a fixed level, Transport+ACI without tail shrinkage, and the full scheme. Every other constant is held fixed. We report the marginal contribution of each component to the interval score with a bootstrap interval over whole days, and, as the reviewer explicitly asked, the runtime each component costs.}

Over the whole test period:

\begin{center}
\begin{tabular}{@{}lccc@{}}
\toprule
Component & $\Delta$ interval score & 95\% CI & Runtime \\
\midrule
Online adaptation, over the static split & $-37.8$ & $[-60,-10]$ & $+0.2$~s \\
Transport alone, over the static split & $-37.8$ & $[-56,-19]$ & $+7.1$~s \\
Transport added on top of ACI & $\mathbf{+74.2}$ & $[+47,+105]$ & $+0.4$~s \\
Tail shrinkage added on top of transport & $-6.5$ & $[-10,-3]$ & $+6.1$~s \\
Full pipeline, over the static split & $+32.4$ & $[-5,+75]$ & $+6.7$~s \\
\bottomrule
\end{tabular}
\end{center}

The answer to the reviewer's question, whether the newly introduced mechanisms and their computational complexity bring tangible improvements, is no for two of the three.

Online adaptation earns its cost: $-215.8$~MWh in the transition window, $[-274,-161]$, for two tenths of a second. Score transport does not, once adaptation is present: it \emph{worsens} the interval score by 74.2~MWh over the whole test with an interval that excludes zero, and it worsens it in four of the five windows separately. Transport without adaptation does improve on the static split, by an amount indistinguishable from what adaptation alone achieves, at thirty-five times the runtime; the two are substitutes and the cheaper one is at least as good. The tail shrinkage, which is the most intricate device in the construction and the one the reviewer singled out, has an effect of $-6.5$~MWh with an interval of $[-10,-3]$ while consuming 6.1 of the 6.7 seconds the full pipeline costs, roughly ninety per cent of its computational budget. The interval excludes zero, so the device does something, but it is six tenths of one per cent of an interval score of about 1{,}030 for ninety per cent of the runtime, and we report it as not paying for itself rather than as undetectable.

Taken together the full pipeline gains 176.3~MWh in the transition window and loses 100.0 and 104.2 in 2025-S1 and 2025-S2, all three significant, netting $+32.4$ over the whole test with an interval containing zero. Over the evaluation period as a whole it is not better than the static split it replaces.

\where{Section~5.5 is a new subsection. Table~8 gives the marginal contributions with intervals and runtimes; Figure~9 shows them window by window. Section~4.6 now states which parts of the construction are new (see R2.4) and Section~4.6 closes by noting that neither survives its ablation, which we judged better placed there than left for the reader to discover in Section~5.}

\subsection*{R1.4 --- Hyperparameter selection and the reservation of the test set}

\begin{revcomment}
The rationale for selecting key hyperparameters and the model tuning procedure are not sufficiently validated, which may undermine the objectivity of experimental results. Independent validation or systematic selection procedures are missing for hyperparameters such as window length, gamma, and parameters of the base forecasting model. This may introduce result selection bias, and it cannot be confirmed whether the reported performance depends on a specific favourable parameter combination. It is recommended that major hyperparameters be determined via an independent validation set or rolling origin validation in the training and calibration phases. Sensitivity analysis for key hyperparameters including window length and gamma should be carried out. In addition, the final test set must be strictly reserved for one off performance evaluation to improve the reliability and reproducibility of experiments.
\end{revcomment}

\action{We accept the criticism and have redone the selection. We should be explicit about what the submitted version did: $\gamma\in\{0.02,0.05\}$ and the 60-day window came from a synthetic prototyping stage and from common practice, not from a formal selection on real data, and both values of $\gamma$ were reported in the test table. That does not meet the standard the reviewer sets.}

\textbf{Base-model hyperparameters.} These were already handled correctly and we now say so where a reader can check it. They were fixed a priori (400 trees, learning rate 0.05, depth 6, min-child-weight 5, subsample and colsample 0.9), never tuned against 2024--2026, and the models are fitted only on targets up to 31 December 2023. The dispersion $\sigma$ is estimated out of fold by five-fold cross-fitting inside the training period. The new multi-quantile model uses the identical settings, which is what makes the R1.1 comparison a comparison of model form rather than of tuning effort.

\textbf{Calibration-layer hyperparameters.} Base-model predictions begin on 1 January 2024, so the only material available before the test period is the 244-day calibration window. We split it into an inner calibration pool (January to April 2024) and an inner validation set (May to August 2024), evaluated in rolling origin with the same seven-day refresh and seven-day embargo as the real protocol. The criterion is the mean interval score, which is proper and which takes the value $+\infty$ if any interval is infinite, so a configuration that escapes into uninformative intervals is disqualified by the criterion rather than by a rule we impose. Grids: $\gamma\in\{0.005,0.01,0.02,0.05,0.10,0.20\}$ and window $\in\{30,45,60,90,120\}$ days.

The selection returns $\gamma=0.005$ for pure ACI and $\gamma=0.005$ with a 120-day window for the transported variant. Both are an order of magnitude gentler than what we reported. The difference is not cosmetic: at $\gamma=0.05$ between 34 and 42 per cent of validation intervals are infinite, whereas at the selected values none is.

\textbf{Test set.} It is used once. The $6\times5$ sensitivity analysis the reviewer asks for is reported on the test set, but only as a display after the selection is closed, and the manuscript says so in those words at both places where it appears. It shows that coverage is flat in $\gamma$ (between 89.9 and 90.1 per cent across the whole grid for the transported variant) and that the fraction of infinite intervals is what varies, from zero at $\gamma\le0.01$ to 22 per cent at $\gamma=0.20$.

\where{Section~4.7 is a new subsection giving the protocol and the outcome. Section~5.8 reports the boundary and chronology sensitivity. The selected configuration appears as its own block in Table~9 so that it can be compared with the values of the submitted version on equal terms. The grids, the criterion and the seed are in Appendix~C.}

\subsection*{R1.5 --- Metrics for one-sided limits and the treatment of infinite bounds}

\begin{revcomment}
The current evaluation framework is insufficient for measuring the performance of final prediction intervals generated by different calibration methods. This paper mainly uses empirical coverage, mean interval width computed only over finite intervals, and the fraction of infinite intervals to evaluate conformal calibration methods. Note that CRPS characterises only the base predictive distribution and cannot reflect performance differences among calibration schemes such as Static split, ACI and Transport+ACI. Especially for scenarios that produce infinite upper prediction bounds, calculating mean interval width solely on finite intervals in Table 2 may give an overly optimistic view of sharpness for certain methods. On top of existing coverage and width metrics, the authors should adopt comprehensive metrics suitable for one sided prediction limits, and clarify how infinite prediction bounds are handled in overall performance assessment. This allows a more fair comparison of the trade off between reliability and sharpness across different calibration methods.
\end{revcomment}

The reviewer identified a real defect, and the direction of the bias is exactly as stated: a method that escapes into infinite intervals removes its widest cases from the width average and therefore looks sharper for having failed.

\action{We adopt the interval score of the one-sided prediction limit as the primary metric, we state explicitly how infinite bounds enter the overall assessment, and, going beyond the comment, we evaluate a modification that removes them.}

\textbf{The metric.} For the interval $[0,U]$ at level $1-\alpha$ the scoring rule of Gneiting and Raftery reduces to
\[
\mathrm{IS}_\alpha(U,y)=U+\tfrac{2}{\alpha}\max(y-U,0),
\]
which is proper, is in MWh, penalises width and non-coverage on the same scale, and is $+\infty$ whenever $U$ is. It cannot be improved by escaping, which is precisely the property the mean finite width lacked.

\textbf{The treatment of infinite bounds.} We report three quantities and label them: the interval score restricted to finite intervals, which is comparable across methods but incomplete; the unrestricted mean, which is the verdict and is infinite for any method producing even one infinite limit; and the fraction of infinite intervals, always. Under this convention three configurations of the submitted version, at up to 5.9 per cent infinite, score as infinite, which is the correct verdict for a method returning uninformative intervals.

\textbf{The remedy.} The cause is structural: the adaptive update is allowed to drive $\alpha_t$ to zero or below, at which point the required order statistic does not exist. We evaluated bounding it from below at $\alpha_{\min}\in\{0.005,0.01,0.02\}$. At $\alpha_{\min}=0.005$ the fraction of infinite intervals falls to zero in all five windows and all four adaptive configurations, at a cost of at most half a point of marginal coverage. We recommend the bound as a default. We also note what it does not fix: with the infinities removed and the scores comparable, the transported variants are still worse than pure adaptation, so that finding is not an artefact of the infinities.

\textbf{On the CRPS.} The reviewer is right that it characterises only the base predictive and cannot arbitrate between calibration schemes. It is retained, and the manuscript now says in as many words why it is retained and what it cannot do. The Brier score of the occurrence event is reported on the same footing and with the same caveat.

\where{Section~4.8 defines the metric and the convention for infinite bounds. Section~5.6 reports the full comparison, Table~9. Table~10 reports the effect of the $\alpha$ bound. Every difference between methods now carries a bootstrap interval that resamples whole days, in Tables 8 and 9 and in Figure~9.}

\clearpage
\section{Reviewer 2}

\subsection*{R2.1 --- Literature review}

\begin{revcomment}
The literature review is inadequate. The manuscript contains only ten references and focuses mainly on conformal prediction, while largely omitting curtailment forecasting, BESS operation, transmission congestion, grid capacity, cable ampacity, thermal constraints, and hybrid physics-data modelling. The authors should also cite broader independent literature to maintain balance.
\end{revcomment}

\action{Section~2 has been rewritten and the bibliography expanded from 10 to 43 entries, organised in four subsections that cover each of the areas named. Following the Editor's instruction, there are no grouped citations anywhere in the manuscript: every reference is discussed individually for what it establishes and why it bears on this problem.}

Coverage of the areas the reviewer lists:

\begin{itemize}[leftmargin=1.4em,itemsep=2pt]
\item \textbf{Curtailment forecasting and measurement} (\S2.1): Henriot on the efficient level of curtailment; Joos and Staffell on integration costs in Britain and Germany; Tang et al.\ on the Chinese experience; Schermeyer et al.\ on congestion management as the mechanism; Nycander et al.\ on Nordic curtailment decomposed by constraint; Yildiz et al.\ on metered low-voltage curtailment; Bunodiere and Lee on rule-based curtailment prediction in Kyushu, which is the closest published forecasting analogue.
\item \textbf{BESS operation} (\S2.2): Denholm and Mai on the storage duration needed to absorb a given curtailment level; Wang et al.\ on unit commitment with storage; Rayit et al.\ on the storage owner's economics of curtailed wind; Mohamad et al.\ on siting storage to minimise curtailment under power-flow constraints.
\item \textbf{Transmission congestion} (\S2.2): Vargas et al.\ on curtailment and storage as congestion-management instruments with ramp limits, on a corridor-constrained system; Fang et al.\ on locational marginal prices under joint renewable and demand uncertainty.
\item \textbf{Grid capacity, cable ampacity and thermal constraints} (\S2.2): Foss et al.\ for the original dynamic ampacity algorithm; Teh and Cotton on the reliability consequences of dynamic thermal rating in wind-integrated networks; Lai and Teh's review; Glaum and Hofmann on the capacity dynamic line rating would unlock on the German grid; Su et al.\ on forecasting transmission capacity itself; Maximov et al.\ on cable ampacity under non-uniform soil temperature; Ocło\'n on soil thermal conductivity against material cost.
\item \textbf{Hybrid physics-data modelling} (\S2.3): Lei et al.\ on physics-informed optimal power flow; Kruse et al.\ on grid frequency; Kirchner-Bossi et al.\ on a hybrid wind-power forecaster. This subsection also states plainly why that route is closed to us: the physics governing curtailment is the operator's security-constrained dispatch, and neither the network model nor the constraint set is public.
\item \textbf{Probabilistic forecasting} (\S2.4), which the reviewer did not list but which the paper needed: Maciejowska et al., Uniejewski and Weron, van der Meer et al., Mangalova and Shesterneva.
\item \textbf{Conformal prediction} (\S2.4), now including Romano et al.\ on conformalized quantile regression, Guan on localised conformal prediction, and the two power-system applications by Hu et al.\ and Jensen et al.
\end{itemize}

On balance: the added references are almost entirely from groups with no connection to ours, across seven countries, and several of them, in particular Bunodiere and Lee, Mangalova and Shesterneva, and Romano et al., report methods that outperform or substitute for our own. Romano et al.'s construction is the benchmark that turns out to dominate our pipeline in Section~5.6, and we say so there.

\textbf{On verification.} Every DOI in the bibliography was resolved against Crossref and checked for title, authors, journal and year; the script that does it and the resulting table are in the companion repository. Three candidate references were dropped because their DOIs did not resolve cleanly, and none was added on the strength of a recollection. We would rather submit a shorter bibliography than one we cannot verify.

\where{Section~2, rewritten in four subsections. Bibliography, 43 entries.}

\subsection*{R2.2 --- Evidence for a storage-driven regime change}

\begin{revcomment}
The BESS-related regime change is not sufficiently demonstrated. The authors should provide BESS capacities, commissioning dates, locations, charging-discharging data, and their relationship to affected renewable plants. Changes in renewable capacity, transmission outages, demand, hydrology, and marginal prices should also be evaluated. Otherwise, the wording should be changed from ``storage-driven regime change'' to ``a regime change coincident with BESS deployment.''
\end{revcomment}

\action{We agree with the reviewer's standard of evidence and cannot meet it with this dataset, so we take the second branch of the alternative --- and in fact go further than the wording proposed.}

We want to be clear about what happened here, because we think the reviewer's reading was reasonable and the fault is ours. The submitted version already made no causal attribution to storage. It dated the distributional change by a formal change-point analysis rather than by any technology's entry, defined the evaluation window as the one of maximum divergence from the calibration distribution, and confined storage to a single bounded hypothesis in the discussion, stated alongside the upper-bound nature of the supporting energy balance and three alternative mechanisms. Our text was evidently not clear enough that this was so, and the reviewer's comment let us find why.

A key in our code, \texttt{test\_ramp}, had survived from an earlier version of the analysis into the published table, carried by a footnote that mentioned storage deployment. A reader encountering a column named ``ramp'' with a footnote about storage would reasonably reconstruct a causal thesis, and that is what happened. It is a nomenclature residue and we should have caught it.

What we did:

\begin{enumerate}[leftmargin=1.6em,itemsep=2pt]
\item Renamed \texttt{test\_ramp} to \texttt{test\_transition} across the codebase, the versioned result files, the tables and the burned-in figure labels. We then verified that the rename changed labels only: the regenerated result files are identical to the submitted ones in every numerical cell, which we checked by direct comparison against a frozen copy.
\item Removed the footnote mentioning storage deployment from what is now Table~3.
\item Audited every remaining mention of storage, batteries and BESS in the manuscript. Storage now appears in three places: in Section~2.2 as literature, in Section~6.3 as one bounded hypothesis among five candidate mechanisms for a distinct phenomenon (the trend break at the turn of 2024 to 2025, not the distributional shift), and in the data-availability statement describing the supplementary registry.
\item Annotated the dated internal documents that use the old vocabulary, rather than rewriting them, since they are a record of what was thought at the time.
\end{enumerate}

On the evidence the reviewer asks for, we state plainly in Section~6.3 what we can and cannot do. We have compiled from the operator's public records a registry of the storage systems of the system with their commercial operation dates and locations, which is supplied as supplementary material. We do not have charge-discharge data, network topology, transmission outage records or nodal prices, and without them the analysis the reviewer describes cannot be done. The literature we now cite makes the point sharper: Mohamad et al.\ show the effect of storage on curtailment depends on its position relative to the binding constraint, and Denholm and Mai show it depends on duration as much as on power, so a capacity-and-date registry alone is not sufficient even in principle.

We did not construct a causal study of storage, and we want to say why explicitly rather than leave it implicit. The methodological contribution never depended on knowing the cause of the shift. A calibration layer that only works when the analyst can name the cause would be of little operational use, and the revision makes that explicit rather than papering over it. The evaluation window is now justified on a measurable criterion, maximum divergence from the calibration distribution, which is the property that actually matters for a conformal layer, and Section~5.8 shows the conclusions are unchanged across six alternative boundary definitions.

\where{Section~4.8 and Table~3, footnote removed and column renamed. Sections~3.4, 5.2 and 6.3. Figures~5 and 6, labels regenerated. Section~7, limitation on causal identification. Code and result files throughout the companion repository.}

\subsection*{R2.3 --- Chronology and change-point sensitivity}

\begin{revcomment}
The regime timeline is inconsistent. The manuscript refers to a 2025 change, while the evaluation defines October-December 2024 as the BESS ramp. The authors should provide one consistent chronology and conduct sensitivity tests using alternative change-point dates.
\end{revcomment}

The inconsistency was real and the reviewer is right to insist on a single chronology.

\action{There is now one chronology, stated in Section~3.4 and used everywhere without variation, and it is regenerated by a versioned script rather than described. The reference to a 2025 change has been removed: it does not survive the analysis.}

The chronology: the distributional shift is a ramp of roughly eighteen months, with change points in May 2023 ($+1.05$ in mean log-magnitude) and January 2024 ($+0.76$), detected by both PELT and binary segmentation under an $L_2$ cost and a BIC penalty. There is no break in October 2024 and none in 2025. The change is essentially complete before 2025 begins: the Wasserstein-1 distance between October--December 2024 and the same quarter of 2025 is 0.085, below the 0.115 threshold of a whole-day permutation null and therefore not distinguishable from no change.

\textbf{Sensitivity to alternative change-point dates}, in two forms.

First, on the detection itself. We repeated the analysis over 180 configurations, crossing two strata, three deseasonalisation methods (calendar-month median, calendar-month mean, robust STL), five penalty multipliers and three minimum segment lengths, with both algorithms. \emph{No configuration returns a break in 2025, and none returns a break in October 2024.} The two dates that do appear are the most frequent across the whole grid.

Second, on the consequences for the evaluation. We re-ran the calibration comparison on six alternative definitions of the transition window, from two months to seven, whose divergence from the calibration distribution ranges from 0.99 to 1.34. The ordering of the methods is unchanged in every case, and the size of the apparent adaptive gain tracks the divergence, which is the same relationship established in Section~5.3.

We also record a correction of process. In preparing this revision we found that our internal log described the change-point penalty as $3\hat\sigma^2\log n$, which does not reproduce the published dates; the standard BIC penalty $\hat\sigma^2\log n$ does, returning both dates and both jump sizes exactly. The published result was correct; the internal description of how it was obtained was not. The penalty is now stated in the manuscript and the sensitivity grid spans it.

\where{Section~3.4, the single chronology with the 180-configuration robustness check. Section~5.8 and Table~12, the boundary sensitivity. Figure~3, regenerated under the stated penalty and returning the same two change points. Appendix~C, the penalty and the grids.}

\subsection*{R2.4 --- Methodological novelty and reproducibility}

\begin{revcomment}
Methodological novelty and reproducibility are unclear. The authors should identify precisely what is new in the PIT, transport, shrinkage, and ACI combination. A complete algorithm, parameter settings, calibration-window length, bootstrap procedure, random seed, and ablation study are required.
\end{revcomment}

\action{Section~4.6 now attributes each ingredient explicitly, and states what is new. The algorithm, parameters, window length, bootstrap procedure and seeds are covered under R1.2; the ablation study under R1.3.

On reproducibility we have added two things the comment implies but that the previous version did not have. First, a second permanent assertion in the experiment script. The existing one checks that the hurdle arm reproduces the submitted table cell for cell, which guards the protocol but not the new code: the quantile predictive is a different class, with its own grid interpolation and its own handling of the atom at zero, and an error there would leave the hurdle assertion untouched while silently producing the paper's headline. The new assertion builds a quantile predictive from the hurdle's own quantile function and checks that the calibration layer over it reproduces the validated arm; it agrees to 0.10 percentage points of coverage and 0.59 per cent of width, which is the expected discretisation error of a 54-level grid, and it runs on every execution rather than only for one arm.

Second, the experiment was reproduced independently from a clean clone, in an environment resolved from the specification file rather than copied, in which four of the five principal libraries resolved to versions other than the reference ones. The retrained multi-quantile base model matched the versioned one by SHA-256 checksum and the regenerated tables matched cell for cell. We have since pinned the library versions, because that agreement depended on a solver resolving similarly and will not hold indefinitely.}

\textbf{Taken from the literature, unmodified:} the probability-integral-transform score is the distributional-conformal score of Chernozhukov et al.; the randomisation of the atom for a mixed distribution is the classical device of Smith; the online update of the miscoverage level is the adaptive conformal inference of Gibbs and Cand\`es; the monotone one-dimensional transport map is the increasing rearrangement, standard in optimal transport and, for quantile curves, the rearrangement of Chernozhukov et al.; bagging a bootstrap estimate and shrinking towards a stable reference in proportion to its variance is ordinary shrinkage practice.

\textbf{New here:} the combination, which is applying a transport of the \emph{score} distribution, rather than of the response or the covariates, as a sharpness-only step inside an online scheme that carries the coverage control, so that the data-dependence of the map cannot damage validity; and the device, which is the variance-proportional tail shrinkage that interpolates level by level between transporting and not transporting according to the bootstrap dispersion of the estimated map.

We then report, in the same subsection, how each fares under ablation, and the two answers differ. The combination does not survive: with the online update in place, adding the transport worsens the interval score by 74.2~MWh over the whole test, with an interval excluding zero. The device does survive its ablation but not its price: the shrinkage improves the score by 6.5~MWh, with a bootstrap interval that also excludes zero, which is six tenths of one per cent of the score for ninety per cent of the runtime. We regard stating that plainly as more useful than the novelty claim it replaces, and it is what the reframing of the paper is built on.

\where{Section~4.6, attribution and novelty. Algorithm~1 and Figure~4, the complete procedure. Appendix~C, Table~C.13, every parameter and all four seeds (42 for base-model training, 20260720 for the PIT atom randomisation and the transport bootstrap, 11 for the CRPS sampling, 20260901 for the day-block bootstrap and permutation nulls). Appendix~D, environment and commands. Section~5.5, the ablations.}

\subsection*{R2.5 --- Conformal validity under panel dependence}

\begin{revcomment}
The conformal validity under panel dependence requires stronger treatment. Plant observations are strongly correlated by date. Clustered standard errors quantify uncertainty in reported coverage but do not restore exchangeability. Block-, group-, or plant-level conformal calibration should be tested, with coverage reported by technology, plant size, region, and high-curtailment events.
\end{revcomment}

The reviewer is correct, and the distinction drawn here is exactly right: a clustered standard error widens an interval around a reported coverage figure, but the finite-sample guarantee rests on exchangeability, not on the width of that interval. We had conflated the two.

\action{We quantify the dependence, test all three calibration schemes the reviewer names, and report coverage disaggregated on all four axes.}

\textbf{Quantification.} On the calibration set, the intraclass correlation of the PIT score by date is 0.238. With 108.8 plants per day this gives a design effect of 26.7, so the effective calibration sample is about 1{,}000 rows rather than the 26{,}552 at which the finite-sample guarantee was nominally stated. The manuscript now says this where the guarantee is claimed, not only in a caveat.

\textbf{Block conformal.} Subsampling one plant per date makes the calibration scores approximately independent, at the cost of reducing them to 244. Averaged over 200 replicates with a fixed seed, the conformal quantile moves from 0.9242 to 0.9263, while the replicate standard deviation is 0.0112. The point estimate is close to unbiased; what the pooled construction misstates is the precision, not the level. This is a correction to the interpretation rather than an improvement to the method, and we present it as such.

\textbf{Group conformal.} Mondrian calibration by technology and by capacity tercile changes little. By region it is \emph{harmful}: the worst regional deviation from nominal rises from 4.8 to 7.2 percentage points, because several regions supply too few calibration days for a stable quantile. Plant-level calibration has the same problem in stronger form, with the same 7.2-point worst regional deviation despite the best marginal interval score of the group. We report this as a negative result: partitioning the calibration set does not buy conditional validity for free when the partition is fine relative to the effective sample size.

\textbf{Disaggregated coverage}, all four axes, for every method. The result that most deserves emphasis runs against our own method. On high-curtailment days, the 254 test days above the calibration-window ninetieth percentile of system-level curtailment, the static split covers at 91.9 per cent and Transport+ACI at 86.7 per cent. The economically important days are where the adaptive scheme under-covers most. Plain ACI, at 88.3 per cent, sits between, and it has the best conditional coverage overall of any scheme tested, with a worst deviation of 4.7 points, better than every Mondrian variant.

\where{Section~4.3 is a new subsection separating the three guarantee regimes and giving the ICC and design effect. Section~5.7 reports the remedies and the disaggregated coverage, Table~11. Section~7 states as a limitation that none of the three schemes repairs the problem. Appendix~B now reads the static guarantee at the effective rather than the nominal sample size.}

\subsection*{R2.6 --- Benchmarks and metrics}

\begin{revcomment}
The benchmarking and evaluation should be expanded. Add probabilistic baselines such as empirical quantiles, conformalized quantile regression, rolling historical distributions, and zero-inflated parametric models. Report interval score, CRPS, Brier score, conditional coverage, infinite-interval frequency, and uncertainty in differences between methods.
\end{revcomment}

\action{All four benchmarks implemented and all six quantities reported.}

The benchmarks, each under the same nominal level, the same seven-day embargo and the same evaluation rows: a per-plant empirical quantile over a trailing year; a per-plant rolling historical distribution over 60 days refreshed weekly; a one-sided conformalized quantile regression on the multi-quantile predictive, following Romano et al.; and a per-plant zero-inflated lognormal fitted on the calibration window without covariates.

They changed the paper's assessment of its own method. Three of the four attain a better interval score than every scheme built on the hurdle, with day-block bootstrap intervals excluding zero. Two of them buy that with coverage, at 86.5 and 87.2 per cent, and we say so, since an operator needing the stated level cannot use them. The conformalized quantile regression is the exception and is the most informative row of the table: 89.8 per cent coverage, within half a point of nominal, at 276~MWh mean width and an interval score of 601, against 596~MWh and 1030 for the pipeline we proposed. A standard construction from 2019, on a better base predictive, dominates it.

The six quantities: interval score of the one-sided limit, in three declared forms (finite-restricted, unrestricted, and the infinite fraction); CRPS of the base predictive, with an explicit statement of what it cannot arbitrate; Brier score of $\{Y>0\}$, with the caveat that the multi-quantile predictive's occurrence probability is read off a 0.02-resolution grid rather than from a dedicated classifier; conditional coverage on four axes (R2.5); infinite-interval frequency, everywhere; and uncertainty in every difference between methods, by a bootstrap resampling whole days with 2{,}000 replicates.

\where{Section~4.8, the metric set. Section~5.6 and Table~9, the benchmarks and full comparison. Table~10, the $\alpha$ bound. Section~5.7 and Table~11, conditional coverage. Tables~8 and 9 and Figure~9 carry the bootstrap intervals.}

\subsection*{R2.7 --- Moderation of the claims}

\begin{revcomment}
The claims should be moderated. The conclusion that point forecasting offers little value applies only to the tested models, features, horizon, and MAE metric. The Abstract and Conclusion should clearly distinguish finite-sample validity under exchangeability, long-run adaptive control, and empirical coverage in the dependent real dataset.
\end{revcomment}

\action{Both requests implemented literally.}

\textbf{The point-forecasting claim} is now scoped in the same paragraph in which it is made, in the reviewer's own terms: it holds for the five model classes tested, the twelve-feature set, the seven-day horizon and mean absolute error as the criterion. The manuscript adds that a model with access to the operator's network model, weather forecasts, unit-commitment schedules or storage state of charge would be solving a better-posed problem, and that a criterion weighting large-magnitude days might order the same models differently. The abstract no longer asserts the general claim, and the conclusion states it as ``within the models, features, horizon and metric tested''.

\textbf{The three guarantees} are now separated explicitly wherever they appear:

\begin{itemize}[leftmargin=1.4em,itemsep=2pt]
\item \emph{Finite-sample marginal validity} of the static split, which requires exchangeability, is an average over draws of the whole calibration set and the test point, and which this panel violates both across the regime boundary and within it, at an effective sample size twenty-seven times smaller than the nominal one.
\item \emph{Long-run adaptive coverage control}, which holds for any sequence with no distributional assumption, says nothing at finite time, and is bought by giving up the finite-sample guarantee.
\item \emph{Empirical coverage in a dependent panel}, which is what we measure, is described by neither of the above, and varies across technologies, plant sizes, regions and high-curtailment days in the ways reported.
\end{itemize}

The abstract closes on that distinction, Section~4.3 is a new subsection devoted to it, and the conclusion restates it as its penultimate paragraph.

\textbf{Beyond the comment.} Since the reviewer's concern is over-claiming, we note that the central claim of the submitted version, the one-third reduction, is no longer made as a property of the method. It is reported as what it is: a result specific to one base model, and a measurement of that model's miscalibration.

\where{Abstract, final two sentences. Section~4.3, the three regimes. Section~5.1, the scope of the point-forecasting result. Section~7, the limitations section. Section~8, conclusion.}

\clearpage
\section{Structured questions and further changes}

\subsection*{Q3 --- Statistical analyses and reporting (Reviewer 1: No)}

\action{Three changes address this. Hyperparameter selection is now by rolling-origin validation with the test set reserved for a single evaluation (R1.4). The primary metric is a proper scoring rule appropriate to one-sided limits, with an explicit convention for infinite bounds (R1.5). And every difference between methods is reported with a bootstrap interval that resamples whole days rather than rows, which is the unit of dependence in this panel (R2.5). We no longer report a difference without an interval, and we state in the text when an interval contains zero.}

\subsection*{Q5 --- Interpretation and conclusions supported by the data (Reviewer 1: No; Reviewer 2: needs more expansion)}

\action{This is the same problem the R1.1 experiment uncovered, and it is the reason the paper was reframed. The conclusion that the adaptive scheme delivers a one-third improvement in sharpness under distribution shift was not supported by the data: it was supported by one base model. The revised paper draws the conclusion the data do support, which is a diagnostic relationship between base-model over-dispersion and apparent calibration-layer gain, and reports the ablations and benchmarks that constrain it.}

\subsection*{Q4 --- Additional tables or figures (both reviewers: Yes)}

\action{Four figures added: Figure~4, the algorithm flowchart (requested in R1.2); Figure~7, the model-agnosticism result and the diagnostic relationship; Figure~8, the capacity ladder and the diagnostic over the forty cells; Figure~9, the marginal contribution of each component with bootstrap intervals. Ten new tables: the three base models compared, the fifteen-cell diagnostic, the capacity ladder, the forty-cell diagnostic, the ablations, the benchmarks, the $\alpha$ bound, the panel-dependence schemes with their disaggregated coverage, the boundary sensitivity, and the complete hyperparameter specification. No existing figure was removed.}

\subsection*{Q7 --- Limitations (both reviewers: No)}

\action{Section~7 is new and dedicated, with ten items: the scope of the point-forecasting conclusion; the failure of exchangeability and the effective sample size; the fact that the diagnostic rests on forty cells from eight predictives that span only two model families, three hurdle variants and five budget variants of one quantile family, on a single dataset; the boundaries of the multi-quantile predictive, including its quantile ceiling and its coarse occurrence probability; that differences of one to three coverage points lie within the clustered standard error; the shortness of the inner validation set; that the causes of the distributional change are not identified and cannot be with this dataset; the six descriptive statements corrected in this revision; and that the dataset records curtailment only, with no network state, dispatch, price or storage operation.}

\subsection*{Q8 and Q9 --- Structure, flow and language (both reviewers: Yes)}

\action{Section~4 is now subdivided into seven subsections separating the base models, the score, the guarantees, the responses to the shift, the algorithm, the novelty statement, the selection protocol and the evaluation, which were previously interleaved. Section~5 is subdivided into seven result subsections, each answering one question. Discussion and limitations, previously mixed, are now Sections~6 and 7. The manuscript has been edited for language throughout, with attention to the passive constructions and long compound sentences of the submitted version.}

\subsection{Corrections we made on our own initiative}
\label{sec:correcciones}

While regenerating every number from versioned scripts, six descriptive statements in Section~3 of the submitted version did not reproduce. None affects a result; all were descriptions of the dataset, and all are corrected in the revised text.

\begin{center}
\footnotesize
\begin{tabular}{@{}p{0.42\textwidth}p{0.50\textwidth}@{}}
\toprule
Submitted version & Regenerated \\
\midrule
Wasserstein-1 of the non-transition windows against calibration: ``0.28 to 0.76'' & 0.28 to 0.95. The September 2024 window is 0.947 and fell outside the range quoted \\
Permutation threshold for the October--December 2024 against 2025 contrast: 0.105 & 0.115. The verdict, 0.085 and not distinguishable from zero, is unchanged \\
Occurrence for northern solar ``stays within 0.71 and 0.85 throughout 2024 to 2026'' & True of the window averages (0.750 to 0.815); at monthly resolution the range is 0.635 to 0.877. Now stated as a window aggregate \\
``The peak hour of the normalised intraday profile is hour 16 in every half-year'' & Hour 16 through 2025-H1 and hour 17 in the last two half-years; the two hours differ by less than 0.4 percentage points of profile mass \\
Total variation between consecutive half-years ``between 0.04 and 0.06, never above 0.08'' & 0.032 to 0.079 from 2023-H1 onwards \\
The profile centroid ``moves by less than fifteen minutes between half-years'' & Up to 22 minutes \\
\bottomrule
\end{tabular}
\end{center}

Two further matters of process, disclosed for completeness. The change-point penalty was described internally as $3\hat\sigma^2\log n$; the standard BIC penalty $\hat\sigma^2\log n$ is what reproduces the published dates and jump sizes, and it is now stated in the manuscript. And the year-over-year growth figures quoted for the trend break used two different conventions for the pre-break level between the two strata; both conventions are now computed and reported, and the manuscript uses one consistently.

The complete verification, twenty-four numerical claims of which eighteen reproduced unchanged, is a versioned artefact in the companion repository.

\end{document}
